\section{Moduly překladače}

\subsection*{scanner.c + scanner.h}
Implementace scanneru v podobě stavového automatu. Čte se postupně znak po znaku a mění se stav.
Pokud se narazí na neočekávaný znak, tak to způsobí lexikální chybu. Pokud se za určitých podmínek skončí v koncovém stavu,
vrátí se série načtených znaků jako jeden token. Hlavní funkcí modulu je \texttt{get\_token},
která vrací právě jeden token při zavolání, pokud nenastala chyba.

\subsection*{parser.c + parser.h}
Implementace parseru, konkrétně syntaktické analýzy pomocí LL gramatiky a přísluš\-ných přechodů.

\subsection*{ast.c + ast.h}
Implementace abstraktního syntaktického stromu, který se využívá jako hierarchická reprezentace načteného kódu.
Jde o jediný strom skládající se z uzlů s určitými vlastnostmi podle typu.

\subsection*{symtable.c + symtable.h}
Tabulka symbolů je implementována jako zásobník, respektive linked list výškově vyvá\-žených binárních stromů.
Konkrétní typ binárních stromů je AVL strom. Každý strom v linked listu reprezentuje jeden scope deklarací identifikátorů.
Nejníže položený strom obsahuje identifikátory funkcí, ostatní neobsahují identifikátory funkcí, pouze proměnných.

\subsection*{semantic.c + semantic.h}
Sémantická analýza propojuje abstraktní syntaktický strom (ASS) s tabulkou symbolů.
Postupně se zanořuje hlouběji do ASS a sekvenčně kontroluje sémantiku. V případech, kdy narazí na výraz, probíhá kontrola rekurzivně.

\subsection*{codegen.c + codegen.h}
Generování kódu převádí abstraktní syntaktický strom na cílový kód.
